% Proposed alternative Introduction (version 2), shorter and in a plainer style.
% Not yet input in main.tex. Part A replaces the whole Introduction.
% Part B (the Bayes factor as a product of forecasts) moves to the start of
% Section 4.1 "Scores are forecasts", so that \eqref{eq:prequential} keeps working.

%======================================================================
% PART A. Introduction
%======================================================================
\section{Introduction}\label{sec:intro}

A digital twin is useful only as long as it behaves like the system it
copies. Checking this is hard when the system is being steered while we
watch. Consider a multi-arm multi-stage (MAMS) trial of insulin-dosing
strategies \citep{whitehead1997design,magirr2012generalized}. At each interim
analysis some arms are dropped and their participants move to other arms,
taking their glucose history with them. An in-silico simulator of type~1
diabetes exists and is accepted by the US Food and Drug Administration for
preclinical testing \citep{kovatchev2009insilico}. We would like to check,
while the trial runs, whether the real participants still behave like
simulated ones, and to raise an alarm as soon as they do not, with a
guaranteed false-alarm rate.

The same question comes up with multi-armed bandits
\citep{lattimore2020bandit}, response-adaptive clinical designs
\citep{trippa2012bayesian,villar2015multi}, and AI systems that act in
environments shaped by their own past decisions. Three features make it
difficult. The data are dependent, because what is observed next depends on
what has been seen. The null model is a simulator, which we can run but whose
likelihood we cannot write down. And we want to look at the evidence after
every observation. Fixed-sample tests lose their guarantees under repeated
looks, and the corrections developed for adaptive trials
\citep{bowden2017unbiased,magirr2012generalized} assume a pre-specified
adaptation rule. Bayesian updating copes with adaptive sampling but does not,
by itself, control the type~I error under optional stopping.

E-values were designed for this situation
\citep{vovk2021evalues,ramdas2023game,grunwald2024safe}. An e-process is the
running wealth of a gambler who places one fair bet against the null per
observation. By Ville's inequality, the probability that this wealth ever
exceeds $1/\alpha$ is at most $\alpha$, so the experimenter can look after
every observation and stop whenever she likes. The best known e-process is
the Bayes factor. Written one observation at a time, it multiplies the
posterior predictive density of each new observation and divides by its null
density. Each factor has mean one under the null for every prior, and the
product rewards whoever forecasts the data better, the learner or the null
\citep{dawid1984present}. Learning the parameters along the way costs only a
term of order $\log t$ \citep{clarke1990information}.

We cannot use the Bayes factor here, for two reasons. The null density is not
available, since the simulator has latent dynamics and, under adaptive
allocation, its law changes with the path. Nor is there a trusted model of the
real system to forecast with, since that is precisely what is in question.
Our approach keeps the shape of the Bayes factor, one fair bet per
observation placed before the observation is seen, and replaces both of its
ingredients.

To replace the null density we start from a classifier. A classifier trained
to tell real outcomes from simulated ones estimates the ratio of their
densities from samples alone \citep{thomas2022lfire}. Such an estimate is not
a fair bet, but a Monte Carlo test turns it into one
\citep{besag1989generalized}. At each step we run the simulator $M$ times,
score the $M$ simulated outcomes together with the real one, and bet on the
real outcome's share of the total score, multiplied by $M+1$. If the simulator
is right, the real outcome is just one more draw among $M+1$ exchangeable
ones, and the bet is fair for any score. In an adaptive design the only
delicate point is how to simulate. We run the simulator under the action
actually taken and from the current state of the unit. This is enough for
validity under any allocation rule, and it cannot be dropped: a simulator run
without regard to the action produces false alarms with probability one.

To replace the forecast we look at the bet from a different angle. As a
function of the parameters of the score, the bet is a likelihood of a familiar
kind. In a conditional-logit model, one item is chosen from a set with
probability proportional to its score; here the set is the real outcome
together with its $M$ simulations, and the item chosen is the real one. Seen
as a likelihood, the bet can be updated by Bayes' rule, one observation at a
time, and the next bet can be placed with the posterior predictive, exactly as
in the Bayes factor. Standard regret bounds for Bayesian mixtures
\citep{kakade2005online} then guarantee that, on every sample path, the wealth
stays within a $d\log t$ term of what the best score in the class would have
earned in hindsight.

Our contributions are the following.
\begin{enumerate}
\item \emph{Validity under adaptive actions, and why the conditioning is
needed.} The Monte Carlo e-process is a test martingale for every allocation
rule (Theorem~\ref{thm:valid}). Without conditioning on the action, the
expected bet under the null tends to $1+\chi^2$ for the likelihood-ratio score
and false alarms grow exponentially, at a rate of at most $\log K$ per step
with $K$ arms (Proposition~\ref{prop:necessary},
Corollary~\ref{cor:falsealarm}).
\item \emph{Bayesian updating of the score.} Reading the bet as a likelihood
gives exact posterior-predictive bets, updated in one step per observation,
with a regret bound that holds on every sample path
(Theorem~\ref{thm:simlik}). A cheaper approximate posterior keeps validity
exact and loses a bounded amount of power (Proposition~\ref{prop:approx}).
\item \emph{Adaptive actions as a source of power.} Since validity holds for
any allocation rule, the actions can be chosen to collect evidence quickly. In
an audit of a language-model API, choosing the prompts adaptively reduces the
number of queries by up to 61\%, for our scores and for those of concurrent
methods (Section~\ref{sec:exp-llm}).
\item \emph{The power of an arbitrary score.} A score implicitly bets that
reality looks like the simulator reweighted by the score, a tilted version of
the null. Its growth rate is the Kullback--Leibler information against the
simulator minus the distance from reality to this tilted null, and $M$
simulations keep a fraction $M/(M+1)$ of it for local alternatives
(Theorem~\ref{thm:growth}).
\end{enumerate}

The pieces are independent of each other. Any score fixed before the new
outcome is seen gives a valid test, so a network trained in advance can be
used in the same adaptive design. Our update rules are what we recommend when
nothing is known about the alternative, which is the usual situation when a
simulator is validated: a twin tends to fail in ways its designers did not
anticipate, and a score tuned to a guessed departure has no power against
another.

In the insulin trial, with five strategies over three stages, only the monitor
that runs the twin from each participant's state under the assigned arm
controls false alarms (2\% of trials, against 95\% for a twin that ignores the
participants' history). One e-process per arm, combined by e-BH
\citep{xu2021unified}, identifies the arm on which the twin fails.

Three recent papers use the same Monte Carlo bet for independent data
\citep{dombowsky2026besag,forre2026anytime,cortinovis2026slice}; for a single
test it is the soft-rank e-value of \citet{wang2022false}. We developed our
construction independently (Appendix~\ref{app:timeline}). Their main
developments assume independent observations and a simulator without state.
\citet{forre2026anytime} list dependent data as an open problem, and
\citet[Appendix~C.1.1]{cortinovis2026slice} note that the construction extends
under conditional exchangeability and prove validity given independent
covariates. They also learn the score differently: \citet{forre2026anytime}
refit it by maximising its empirical growth, and \citet{cortinovis2026slice}
train a neural ratio estimator before the data arrive. Section~\ref{sec:related}
lists in detail which of our results coincide with theirs.

Section~\ref{sec:method} gives the method and its validity,
Section~\ref{sec:power} its power, and Section~\ref{sec:predictive} the
predictive view and the update rules (Algorithm~1). Section~\ref{sec:experiments}
contains simulations, the comparison with concurrent methods and the
language-model audit, and Section~\ref{sec:trial} the trial.
Section~\ref{sec:discussion} discusses limitations.

%======================================================================
% PART B. Opening of Section 4.1 "Scores are forecasts"
% (the Bayes factor material moved from the Introduction)
%======================================================================
% Insert right after \subsection{Scores are forecasts}\label{sec:scores}

The Bayes factor is the model for what follows. For a composite alternative
with model $q_\theta$ and prior $\pi_0$, it is the ratio of the marginal
likelihood $m(x_{1:t}) = \int \prod_{i \le t} q_\theta(x_i \mid x_{1:i-1})\,
\pi_0(d\theta)$ to the null density. By the chain rule,
\begin{equation}\label{eq:prequential}
\mathrm{BF}_t = \prod_{i=1}^{t} \frac{m_i(x_i)}{p_0(x_i \mid x_{1:i-1})},
\qquad
m_i(x) = \int q_\theta(x \mid x_{1:i-1})\,\pi_{i-1}(d\theta),
\end{equation}
where $\pi_{i-1}$ is the posterior after $i-1$ observations, so $m_i$ is the
posterior predictive of $x_i$. (For a coin with a uniform prior,
$m_i(\text{heads}) = (k+1)/(i+1)$ after $k$ heads in $i-1$ tosses.) Each
factor has conditional mean $\int m_i = 1$ under the null, whatever the prior;
$\log \mathrm{BF}_t$ is the difference between the cumulative log scores of
the learner and of the null \citep{dawid1984present}; and for i.i.d.\ data
and a regular model with $d$ parameters
\[
\log \mathrm{BF}_t = \max_\theta \sum_{i \le t} \log\frac{q_\theta(x_i)}{p_0(x_i)}
- \frac d2 \log t + O_P(1)
\]
\citep{clarke1990information}. The Monte Carlo e-process has the same
structure, with the simulations in place of $p_0$ and the score in place of
$q_\theta$; the rest of this section makes the correspondence precise.
